% 对应题号 M1-46，复用已核对的本题几何图。
% 修订：启用 xeCJK 自动换行；保留原题参数与物理关系。
\documentclass[tikz,border=14pt]{standalone}
\usepackage{fontspec}
\usepackage{xeCJK}
\usepackage{amsmath,amssymb}
\usepackage{xcolor}
\usepackage{tikz}
\usetikzlibrary{arrows.meta,calc,patterns,positioning,decorations.pathmorphing}
\IfFontExistsTF{Heiti SC}{\IfFontExistsTF{Heiti SC}{\setmainfont{Heiti SC}}{\setmainfont{FandolHei-Regular}}\IfFontExistsTF{Heiti SC}{\setsansfont{Heiti SC}}{\setsansfont{TeX Gyre Heros}}}{\IfFontExistsTF{Songti SC}{\setmainfont{Songti SC}}{\setmainfont{FandolSong-Regular}}\IfFontExistsTF{Songti SC}{\setsansfont{Songti SC}}{\setsansfont{FandolSong-Regular}}}
\definecolor{bluefield}{RGB}{24,76,255}
\definecolor{redfield}{RGB}{232,49,35}
\definecolor{orangefield}{RGB}{224,116,10}
\definecolor{greenfield}{RGB}{18,132,86}
\definecolor{purplefield}{RGB}{112,72,190}
\definecolor{linegray}{RGB}{60,68,78}
\definecolor{softblue}{RGB}{238,243,255}
\definecolor{softgreen}{RGB}{238,249,244}
\definecolor{softorange}{RGB}{255,247,235}
\definecolor{softpurple}{RGB}{247,242,255}
\tikzset{
  arrow/.style={-{Latex[length=3.1mm,width=2.1mm]},line width=1.1pt},
  thinarr/.style={-{Latex[length=2.5mm,width=1.7mm]},line width=.85pt},
  axis/.style={-{Latex[length=2.5mm,width=1.7mm]},line width=.8pt,draw=linegray},
  guide/.style={draw=linegray!55,dashed,line width=.65pt},
  fieldline/.style={-{Latex[length=2.5mm,width=1.7mm]},line width=1pt,draw=bluefield},
  box/.style={draw=linegray,rounded corners=3pt,fill=white,align=center,inner sep=6pt,line width=.75pt},
  bluebox/.style={box,draw=bluefield,fill=softblue},
  greenbox/.style={box,draw=greenfield,fill=softgreen},
  orangebox/.style={box,draw=orangefield,fill=softorange},
  purplebox/.style={box,draw=purplefield,fill=softpurple}
}
\IfFontExistsTF{Heiti SC}{\setCJKmainfont[AutoFakeBold=1.5]{Heiti SC}}{\setCJKmainfont[AutoFakeBold=1.5]{FandolHei-Regular}}
\IfFontExistsTF{Heiti SC}{\setCJKsansfont{Heiti SC}}{\setCJKsansfont{FandolHei-Regular}}
\begin{document}

\begin{tikzpicture}[x=1cm,y=1cm]
\path[use as bounding box] (-10,-4.6) rectangle (10,4.6);
\fill[white] (-10,-4.6) rectangle (10,4.6);
\node[font=\Large\bfseries,text width=18.4cm,align=center] at (0,4.05) {例：刚体旋转处处有旋，势涡的旋度集中在奇点};
\begin{scope}[shift={(-5.2,-.3)}]
  \foreach \r in {1,1.8,2.6}{\draw[bluefield,line width=.9pt,-{Latex[length=2.2mm]}] (\r,0) arc[start angle=0,end angle=300,radius=\r];}
  \node[redfield,font=\Large] at (0,0) {$\odot$};\node[redfield] at (0,3.05) {$\nabla\times\mathbf v=2\boldsymbol\Omega$（指向纸外）};
  \node at (0,-3.05) {$\mathbf v=\boldsymbol\Omega\times\mathbf r$};
\end{scope}
\begin{scope}[shift={(5.2,-.3)}]
  \fill[redfield] (0,0) circle(.11) node[below right] {轴线奇点};
  \foreach \r in {1,1.8,2.6}{\draw[greenfield,line width=.9pt,-{Latex[length=2.2mm]}] (\r,0) arc[start angle=0,end angle=300,radius=\r];}
  \node[greenbox,text width=5.7cm] at (0,-3.65) {$\mathbf v=\Gamma(2\pi r)^{-1}\mathbf e_\phi$\\$r>0$ 处旋度为零，但环量为 $\Gamma$};
\end{scope}
\end{tikzpicture}
\end{document}
